% GENERATED FILE. Edit canonical lessons, then run npm run generate:books.
% canonical-source-hash: fnv1a64:6adbe21a4617f070
% canonical-lessons: TE-C485-cellimpu, TE-C485-nagadu, TE-C485-baki, TE-C485-marakam, TE-C485-eggottu, TE-R485-review-478-485

\chapter{Paying Up}
\label{ch:te-paying-up}
\hlchaptermodality{485}

\begin{chapteropening}
\textbf{By the end of this chapter:} \emph{I can say a payment, cash, money owed, an exchange rate and to dodge.}

Complete the last lesson of chapter 485: I can say a payment, cash, money owed, an exchange rate and to dodge.
\end{chapteropening}

\section[\texorpdfstring{\te{చెల్లింపు} (cellimpu)}{చెల్లింపు (cellimpu)}]{\te{చెల్లింపు} (cellimpu) — a payment}

\label{lesson:TE-C485-cellimpu}



\begin{quote}
Before the new one: say the Telugu for a concession, then the Telugu for to haggle.
\end{quote}

\subsection*{You'll want to know: \te{చెల్లింపు}}
\textbf{\te{చెల్లింపు}} — \emph{cellimpu} — \textquotedblleft{}a payment\textquotedblright{}.

\begin{quote}
\te{బిల్లు} \te{చెల్లింపు} \te{కోసం} \te{కౌంటరుకు} \te{వెళ్ళాను}.
\end{quote}

\begin{quote}
\emph{billu cellimpu kōsaṁ kauṇṭaruku veḷḷānu.}
\end{quote}

\begin{quote}
I went to the counter to pay the bill.
\end{quote}

The noun of \te{చెల్లించు}, to pay. \te{ఆన్లైన్} \te{చెల్లింపు}, an online payment.

\begin{cousinweb}
\textbf{\te{చెల్లింపు}} (\emph{cellimpu}) is native Dravidian, the noun of \te{చెల్లించు}, \textquotedblleft{}to pay\textquotedblright{}, from \te{చెల్లు}, \textquotedblleft{}to pass, to be current, to be paid\textquotedblright{} (DEDR 2781).
\end{cousinweb}

\begin{grammarlens}[title={What you've built}]
One more word for getting around.
\end{grammarlens}

\subsection*{Your turn}
\begin{itemize}
\raggedright
  \item \emph{Say it:} \emph{cellimpu}
  \item \emph{Say it:} \emph{cellimpu}, once more
  \item \emph{Say it:} \emph{billu cellimpu kōsaṁ kauṇṭaruku veḷḷānu.} — I went to the counter to pay the bill.
  \item \emph{Recall:} say the Telugu for a concession, then the Telugu for to haggle, then say \emph{cellimpu} again
\end{itemize}

\begin{tcolorbox}[breakable,colback=teal!4,colframe=teal!35!black,title={Before you move on}]
What does \te{చెల్లింపు} mean? (\textquotedblleft{}A payment\textquotedblright{}.) Say it once more.
\end{tcolorbox}
\section[\texorpdfstring{\te{నగదు} (nagadu)}{నగదు (nagadu)}]{\te{నగదు} (nagadu) — cash}

\label{lesson:TE-C485-nagadu}



\begin{quote}
Before the new one: say the Telugu for to haggle, then the Telugu for a payment.
\end{quote}

\subsection*{You'll want to know: \te{నగదు}}
\textbf{\te{నగదు}} — \emph{nagadu} — \textquotedblleft{}cash\textquotedblright{}.

\begin{quote}
\te{ఆ} \te{దుకాణంలో} \te{నగదు} \te{మాత్రమే} \te{తీసుకుంటారు}.
\end{quote}

\begin{quote}
\emph{ā dukāṇaṁlō nagadu mātramē tīsukuṇṭāru.}
\end{quote}

\begin{quote}
That shop takes cash only.
\end{quote}

\te{నగదు} \te{రూపంలో}, in cash.

\begin{cousinweb}
\textbf{\te{నగదు}} (\emph{nagadu}) is a loan from Arabic \emph{naqd}, \textquotedblleft{}ready money, cash\textquotedblright{}, which came into Telugu through Persian and Hindustani.
\end{cousinweb}

\begin{grammarlens}[title={What you've built}]
One more word for getting around.
\end{grammarlens}

\subsection*{Your turn}
\begin{itemize}
\raggedright
  \item \emph{Say it:} \emph{nagadu}
  \item \emph{Say it:} \emph{nagadu}, once more
  \item \emph{Say it:} \emph{ā dukāṇaṁlō nagadu mātramē tīsukuṇṭāru.} — That shop takes cash only.
  \item \emph{Recall:} say the Telugu for to haggle, then the Telugu for a payment, then say \emph{nagadu} again
\end{itemize}

\begin{tcolorbox}[breakable,colback=teal!4,colframe=teal!35!black,title={Before you move on}]
What does \te{నగదు} mean? (\textquotedblleft{}Cash\textquotedblright{}.) Say it once more.
\end{tcolorbox}
\section[\texorpdfstring{\te{బాకీ} (bākī)}{బాకీ (bākī)}]{\te{బాకీ} (bākī) — money owed; a balance due}

\label{lesson:TE-C485-baki}



\begin{quote}
Before the new one: say the Telugu for a payment, then the Telugu for cash.
\end{quote}

\subsection*{You'll want to know: \te{బాకీ}}
\textbf{\te{బాకీ}} — \emph{bākī} — \textquotedblleft{}money owed; a balance due\textquotedblright{}.

\begin{quote}
\te{కిరాణా} \te{కొట్టుకు} \te{ఇంకా} \te{వంద} \te{రూపాయలు} \te{బాకీ} \te{ఉన్నాను}.
\end{quote}

\begin{quote}
\emph{kirāṇā koṭṭuku iṅkā vanda rūpāyalu bākī unnānu.}
\end{quote}

\begin{quote}
I still owe the grocer a hundred rupees.
\end{quote}

\te{బాకీ} \te{ఉండు}, to owe; \te{బాకీ} \te{తీర్చు}, to settle a debt.

\begin{cousinweb}
\textbf{\te{బాకీ}} (\emph{bākī}) is a loan from Arabic \emph{bāqī}, \textquotedblleft{}remaining, the balance\textquotedblright{}, through Persian and Hindustani.
\end{cousinweb}

\begin{grammarlens}[title={What you've built}]
One more word for getting around.
\end{grammarlens}

\subsection*{Your turn}
\begin{itemize}
\raggedright
  \item \emph{Say it:} \emph{bākī}
  \item \emph{Say it:} \emph{bākī}, once more
  \item \emph{Say it:} \emph{kirāṇā koṭṭuku iṅkā vanda rūpāyalu bākī unnānu.} — I still owe the grocer a hundred rupees.
  \item \emph{Recall:} say the Telugu for a payment, then the Telugu for cash, then say \emph{bākī} again
\end{itemize}

\begin{tcolorbox}[breakable,colback=teal!4,colframe=teal!35!black,title={Before you move on}]
What does \te{బాకీ} mean? (\textquotedblleft{}Money owed; a balance due\textquotedblright{}.) Say it once more.
\end{tcolorbox}
\section[\texorpdfstring{\te{మారకం} (mārakaṁ)}{మారకం (mārakaṁ)}]{\te{మారకం} (mārakaṁ) — an exchange rate}

\label{lesson:TE-C485-marakam}



\begin{quote}
Before the new one: say the Telugu for cash, then the Telugu for money owed.
\end{quote}

\subsection*{You'll want to know: \te{మారకం}}
\textbf{\te{మారకం}} — \emph{mārakaṁ} — \textquotedblleft{}an exchange rate\textquotedblright{}.

\begin{quote}
\te{ఇవాళ} \te{డాలరు} \te{మారకం} \te{రేటు} \te{పెరిగింది}.
\end{quote}

\begin{quote}
\emph{ivāḷa ḍālaru mārakaṁ rēṭu perigindi.}
\end{quote}

\begin{quote}
Today the dollar exchange rate went up.
\end{quote}

\te{మారకం} \te{రేటు}, the exchange rate. \te{మారు}, to change, is in the word.

\begin{cousinweb}
\textbf{\te{మారకం}} (\emph{mārakaṁ}) is native Dravidian, built on \te{మారు}, \textquotedblleft{}to change, to be exchanged\textquotedblright{} (DEDR 4834).
\end{cousinweb}

\begin{grammarlens}[title={What you've built}]
One more word for getting around.
\end{grammarlens}

\subsection*{Your turn}
\begin{itemize}
\raggedright
  \item \emph{Say it:} \emph{mārakaṁ}
  \item \emph{Say it:} \emph{mārakaṁ}, once more
  \item \emph{Say it:} \emph{ivāḷa ḍālaru mārakaṁ rēṭu perigindi.} — Today the dollar exchange rate went up.
  \item \emph{Recall:} say the Telugu for cash, then the Telugu for money owed, then say \emph{mārakaṁ} again
\end{itemize}

\begin{tcolorbox}[breakable,colback=teal!4,colframe=teal!35!black,title={Before you move on}]
What does \te{మారకం} mean? (\textquotedblleft{}An exchange rate\textquotedblright{}.) Say it once more.
\end{tcolorbox}
\section[\texorpdfstring{\te{ఎగ్గొట్టు} (eggoṭṭu)}{ఎగ్గొట్టు (eggoṭṭu)}]{\te{ఎగ్గొట్టు} (eggoṭṭu) — to dodge; to skip paying}

\label{lesson:TE-C485-eggottu}



\begin{quote}
Before the new one: say the Telugu for money owed, then the Telugu for an exchange rate.
\end{quote}

\subsection*{You'll want to know: \te{ఎగ్గొట్టు}}
\textbf{\te{ఎగ్గొట్టు}} — \emph{eggoṭṭu} — \textquotedblleft{}to dodge; to skip paying\textquotedblright{}.

\begin{quote}
\te{అతను} \te{అద్దె} \te{ఎగ్గొట్టి} \te{వెళ్ళిపోయాడు}.
\end{quote}

\begin{quote}
\emph{atanu adde eggoṭṭi veḷḷipōyāḍu.}
\end{quote}

\begin{quote}
He went off without paying the rent.
\end{quote}

Past (I …): \te{ఎగ్గొట్టాను}, \emph{eggoṭṭānu}.

A colloquial word: to dodge a debt, a class or a duty.

\begin{cousinweb}
\textbf{\te{ఎగ్గొట్టు}} (\emph{eggoṭṭu}) is a Telugu compound: -\te{గొట్టు}, from the native \te{కొట్టు}, \textquotedblleft{}to strike\textquotedblright{} (DEDR 2063), and a first part whose origin is not settled.
\end{cousinweb}

\begin{grammarlens}[title={What you've built}]
One more word for getting around.
\end{grammarlens}

\subsection*{Your turn}
\begin{itemize}
\raggedright
  \item \emph{Say it:} \emph{eggoṭṭu}
  \item \emph{Say it:} \emph{eggoṭṭu}, once more
  \item \emph{Say it:} \emph{atanu adde eggoṭṭi veḷḷipōyāḍu.} — He went off without paying the rent.
  \item \emph{Recall:} say the Telugu for money owed, then the Telugu for an exchange rate, then say \emph{eggoṭṭu} again
\end{itemize}

\begin{tcolorbox}[breakable,colback=teal!4,colframe=teal!35!black,title={Before you move on}]
What does \te{ఎగ్గొట్టు} mean? (\textquotedblleft{}To dodge; to skip paying\textquotedblright{}.) Say it once more.
\end{tcolorbox}
\section[(dialogue)]{Words from chapters 478–485 — a review}

\label{lesson:TE-R485-review-478-485}



\begin{quote}
Say the Telugu for a low-pressure area and a gale.
\end{quote}

\subsection*{Your turn}
Say these aloud:

\begin{itemize}
\raggedright
  \item a low-pressure area, a gale, a whirlwind, a rainstorm, to lose strength
  \item a thunderbolt, a torrential rain, to close in, to grow dim, spray
  \item to inundate, to overflow, to snap, to get bogged down, to get stuck
  \item a heatwave, a mirage, a dust storm, to dry up, to crack open
  \item hailstones, ice, snowfall, a breeze, to cool down
  \item to evaporate, to soak in, to blaze, to go out, to put forth leaves
  \item a sale, a consumer, quality, a concession, to haggle
  \item a payment, cash, money owed, an exchange rate, to dodge
\end{itemize}

\begin{tcolorbox}[breakable,colback=teal!4,colframe=teal!35!black,title={Before you move on}]
A low-pressure area? (\textbf{\te{అల్పపీడనం}}, \emph{alpapīḍanaṁ}.) A gale? (\textbf{\te{పెనుగాలి}}, \emph{penugāli}.) A whirlwind? (\textbf{\te{సుడిగాలి}}, \emph{suḍigāli}.) A rainstorm? (\textbf{\te{గాలివాన}}, \emph{gālivāna}.) To lose strength? (\textbf{\te{బలహీనపడు}}, \emph{balahīnapaḍu}.) A thunderbolt? (\textbf{\te{పిడుగు}}, \emph{piḍugu}.) A torrential rain? (\textbf{\te{కుండపోత}}, \emph{kuṇḍapōta}.) To close in? (\textbf{\te{కమ్ము}}, \emph{kammu}.) To grow dim? (\textbf{\te{మసకబారు}}, \emph{masakabāru}.) Spray? (\textbf{\te{తుంపర}}, \emph{tumpara}.) To inundate? (\textbf{\te{ముంచెత్తు}}, \emph{muñcettu}.) To overflow? (\textbf{\te{పొంగిపొర్లు}}, \emph{poṅgiporlu}.) To snap? (\textbf{\te{తెగిపోవు}}, \emph{tegipōvu}.) To get bogged down? (\textbf{\te{కూరుకుపోవు}}, \emph{kūrukupōvu}.) To get stuck? (\textbf{\te{ఇరుక్కుపోవు}}, \emph{irukkupōvu}.) A heatwave? (\textbf{\te{వడగాలి}}, \emph{vaḍagāli}.) A mirage? (\textbf{\te{ఎండమావి}}, \emph{eṇḍamāvi}.) A dust storm? (\textbf{\te{గాలిదుమారం}}, \emph{gālidumāraṁ}.) To dry up? (\textbf{\te{ఎండిపోవు}}, \emph{eṇḍipōvu}.) To crack open? (\textbf{\te{బీటలు} \te{వారు}}, \emph{bīṭalu vāru}.) Hailstones? (\textbf{\te{వడగళ్ళు}}, \emph{vaḍagaḷḷu}.) Ice? (\textbf{\te{మంచుగడ్డ}}, \emph{mañcugaḍḍa}.) Snowfall? (\textbf{\te{హిమపాతం}}, \emph{himapātaṁ}.) A breeze? (\textbf{\te{పిల్లగాలి}}, \emph{pillagāli}.) To cool down? (\textbf{\te{చల్లబడు}}, \emph{callabaḍu}.) To evaporate? (\textbf{\te{ఆవిరవు}}, \emph{āviravu}.) To soak in? (\textbf{\te{ఇంకిపోవు}}, \emph{iṅkipōvu}.) To blaze? (\textbf{\te{మండు}}, \emph{maṇḍu}.) To go out? (\textbf{\te{ఆరిపోవు}}, \emph{āripōvu}.) To put forth leaves? (\textbf{\te{చిగురించు}}, \emph{ciguriñcu}.) A sale? (\textbf{\te{అమ్మకం}}, \emph{ammakaṁ}.) A consumer? (\textbf{\te{వినియోగదారుడు}}, \emph{viniyōgadāruḍu}.) Quality? (\textbf{\te{నాణ్యత}}, \emph{nāṇyata}.) A concession? (\textbf{\te{రాయితీ}}, \emph{rāyitī}.) To haggle? (\textbf{\te{బేరమాడు}}, \emph{bēramāḍu}.) A payment? (\textbf{\te{చెల్లింపు}}, \emph{cellimpu}.) Cash? (\textbf{\te{నగదు}}, \emph{nagadu}.) Money owed? (\textbf{\te{బాకీ}}, \emph{bākī}.) An exchange rate? (\textbf{\te{మారకం}}, \emph{mārakaṁ}.) To dodge? (\textbf{\te{ఎగ్గొట్టు}}, \emph{eggoṭṭu}.)
\end{tcolorbox}
